\chapter{2017年全国II卷（理）}

\section{单选题}

\begin{problem}
\(\dfrac{3+\i}{1+\i}=\)（\quad）

\choices
    {\(1+2\i\)}
    {\(1-2\i\)}
    {\(2+\i\)}
    {\(2-\i\)}
\end{problem}

\begin{answer}
D
\end{answer}

\begin{solution}
分子、分母同乘以 \(1-\i\)，得
\[
\frac{3+\i}{1+\i}=\frac{(3+\i)(1-\i)}{(1+\i)(1-\i)}
=\frac{4-2\i}{2}=2-\i
\]
故选 D.
\end{solution}

\begin{problem}
设集合 \(A=\{1,2,4\}\)，\(B=\{x\mid x^2-4x+m=0\}\). 若 \(A\cap B=\{1\}\)，则 \(B=\)（\quad）

\choices
    {\(\{1,-3\}\)}
    {\(\{1,0\}\)}
    {\(\{1,3\}\)}
    {\(\{1,5\}\)}
\end{problem}

\begin{answer}
C
\end{answer}

\begin{solution}
因为 \(A\cap B=\{1\}\)，所以 \(1\in B\)，代入方程得 \(1-4+m=0\)，从而 \(m=3\). 因此 \(B\) 是方程 \(x^2-4x+3=0\) 的根集，而
\(x^2-4x+3=(x-1)(x-3)\)，
所以 \(B=\{1,3\}\). 故选 C.
\end{solution}

\begin{problem}
我国古代数学名著《算法统宗》中有如下问题：“远望巍巍塔七层，红光点点倍加增，共灯三百八十一，请问尖头几盏灯？”意思是：一座 \(7\text{ 层}\)塔共挂了 \(381\text{ 盏}\)灯，且相邻两层中的下一层灯数是上一层灯数的 \(2\text{ 倍}\)，则塔的顶层共有灯（\quad）

\choices
    {\(1\text{ 盏}\)}
    {\(3\text{ 盏}\)}
    {\(5\text{ 盏}\)}
    {\(9\text{ 盏}\)}
\end{problem}

\begin{answer}
B
\end{answer}

\begin{solution}
设塔顶层有 \(a\) 盏灯，则从顶层到底层各层的灯数构成首项为 \(a\)、公比为 \(2\) 的等比数列. 由总灯数为 \(381\)，得
\(a\frac{2^7-1}{2-1}=381\)，
即 \(127a=381\)，所以 \(a=3\). 故选 B.
\end{solution}

\begin{problem}
如图，网格纸上小正方形的边长为 \(1\)，粗实线画出的是某几何体的三视图，该几何体由一平面将一圆柱截去一部分后所得，则该几何体的体积为（\quad）

\bitmapfigure[max width=0.58\linewidth]{img/2017/national_paper_2_science/q4_fig1.png}

\choices
    {\(90\pi\)}
    {\(63\pi\)}
    {\(42\pi\)}
    {\(36\pi\)}
\end{problem}

\begin{answer}
B
\end{answer}

\begin{solution}
由三视图可知，原圆柱的底面半径为 \(3\)，高为 \(10\)，其体积为 \(\pi\times3^2\times10=90\pi\). 被平面截去的部分是底面半径为 \(3\)、高为 \(6\) 的半个圆柱，其体积为
\(\frac12\pi\times3^2\times6=27\pi\).
所以该几何体的体积为 \(90\pi-27\pi=63\pi\). 故选 B.
\end{solution}

\begin{problem}
设 \(x\)，\(y\) 满足约束条件
\(\begin{cases}
2x+3y-3\le 0,\\
2x-3y+3\ge 0,\\
y+3\ge 0
\end{cases}\)
则 \(z=2x+y\) 的最小值是（\quad）

\choices
    {\(-15\)}
    {\(-9\)}
    {\(1\)}
    {\(9\)}
\end{problem}

\begin{answer}
A
\end{answer}

\begin{solution}
约束条件等价于
\(y\le\frac{3-2x}{3}\)、\(y\le\frac{2x+3}{3}\)、\(y\ge-3\).
三条边界直线的交点分别为 \((-6,-3)\)、\((0,1)\) 和 \((6,-3)\). 因而可行域是以这三个点为顶点的三角形. 在顶点处计算目标函数 \(z=2x+y\)，得到
\(z\bigl(-6,-3\bigr)=-15\)、\(z(0,1)=1\)、\(z(6,-3)=9\).
所以 \(z\) 的最小值为 \(-15\). 故选 A.
\end{solution}

\begin{problem}
安排 \(3\text{ 名}\)志愿者完成 \(4\text{ 项}\)工作，每人至少完成 \(1\text{ 项}\)，每项工作由 \(1\text{ 人}\)完成，则不同的安排方式共有（\quad）

\choices
    {\(12\text{ 种}\)}
    {\(18\text{ 种}\)}
    {\(24\text{ 种}\)}
    {\(36\text{ 种}\)}
\end{problem}

\begin{answer}
D
\end{answer}

\begin{solution}
因为有 \(4\) 项工作、\(3\) 名志愿者且每人至少完成 \(1\) 项，所以工作分配数只能是 \(2\)，\(1\)，\(1\). 先从 \(4\) 项工作中选出由同一人完成的 \(2\) 项，有 \(\mathrm C_4^2\) 种；再把这三组工作分给 \(3\) 名志愿者，有 \(\mathrm A_3^3\) 种. 因而安排方式共有
\(\mathrm C_4^2\mathrm A_3^3=6\times6=36\)
种. 故选 D.
\end{solution}

\begin{problem}
甲，乙，丙，丁四位同学一起去向老师询问成语竞赛的成绩. 老师说：你们四人中有 \(2\) 位优秀，\(2\) 位良好，我现在给甲看乙，丙的成绩，给乙看丙的成绩，给丁看甲的成绩. 看后甲对大家说：我还是不知道我的成绩. 根据以上信息，则（\quad）

\choices
    {乙可以知道四人的成绩}
    {丁可以知道四人的成绩}
    {乙，丁可以知道对方的成绩}
    {乙，丁可以知道自己的成绩}
\end{problem}

\begin{answer}
D
\end{answer}

\begin{solution}
若乙、丙的成绩相同，则甲看到的两人成绩或者都是优秀，或者都是良好. 在这两种情形下，由总人数中恰有 \(2\) 位优秀、\(2\) 位良好，甲都能确定自己的成绩. 但甲说自己仍不知道，所以乙、丙的成绩必不相同.

于是，乙看到丙的成绩后，可以由“乙、丙一优一良”直接确定自己的成绩：丙优秀时乙良好，丙良好时乙优秀. 丁看到甲的成绩后，也可由总人数中两优两良确定自己的成绩：甲优秀时丁良好，甲良好时丁优秀. 但乙只看到丙的成绩，不能据此确定甲、丁二人的具体顺序，故不能确定丁的成绩；同理不能断定乙、丁能知道对方的成绩. 因此只有乙、丁可以知道自己的成绩，故选 D.
\end{solution}

\begin{problem}
执行如图的程序框图，如果输入的 \(a=-1\)，则输出的 \(S=\)（\quad）

\bitmapfigure[max width=0.40\linewidth]{img/2017/national_paper_2_science/q8_fig1.png}

\choices
    {\(2\)}
    {\(3\)}
    {\(4\)}
    {\(5\)}
\end{problem}

\begin{answer}
B
\end{answer}

\begin{solution}
程序在 \(K\le6\) 时循环，每次先将 \(aK\) 加入 \(S\)，再令 \(a=-a\)、\(K=K+1\). 初始 \(S=0\)、\(K=1\)、\(a=-1\)，各次累加的数依次为
\(-1\)、\(2\)、\(-3\)、\(4\)、\(-5\)、\(6\).
因此循环结束时
\[
S=-1+2-3+4-5+6=3
\]
此时 \(K=7>6\)，输出 \(S=3\). 故选 B.
\end{solution}

\begin{problem}
若双曲线 \(C:\dfrac{x^2}{a^2}-\dfrac{y^2}{b^2}=1\)（\(a>0\)，\(b>0\)）的一条渐近线被圆 \((x-2)^2+y^2=4\) 所截得的弦长为 \(2\)，则 \(C\) 的离心率为（\quad）

\choices
    {\(2\)}
    {\(\sqrt3\)}
    {\(\sqrt2\)}
    {\(\dfrac{2\sqrt3}{3}\)}
\end{problem}

\begin{answer}
A
\end{answer}

\begin{solution}
双曲线的一条渐近线可写为 \(bx-ay=0\). 圆的半径为 \(2\)，被渐近线截得的弦长为 \(2\)，所以圆心 \((2,0)\) 到该直线的距离 \(d\) 满足
\(d=\sqrt{2^2-\left(\frac{2}{2}\right)^2}=\sqrt3\).
另一方面，点 \((2,0)\) 到直线 \(bx-ay=0\) 的距离为
\(d=\frac{2b}{\sqrt{a^2+b^2}}\).
于是 \(4b^2=3(a^2+b^2)\)，即 \(b^2=3a^2\). 设双曲线焦半距为 \(c\)，则 \(c^2=a^2+b^2=4a^2\)，从而离心率
\(e=\frac ca=2\).
故选 A.
\end{solution}

\begin{problem}
已知直三棱柱 \(ABC-A_1B_1C_1\) 中，\(\angle ABC=120^\circ\)，\(AB=2\)，\(BC=CC_1=1\)，则异面直线 \(AB_1\) 与 \(BC_1\) 所成角的余弦值为（\quad）

\choices
    {\(\dfrac{\sqrt3}{2}\)}
    {\(\dfrac{\sqrt{15}}{5}\)}
    {\(\dfrac{\sqrt{10}}{5}\)}
    {\(\dfrac{\sqrt3}{3}\)}
\end{problem}

\begin{answer}
C
\end{answer}

\begin{solution}
在底面内取直角坐标，使 \(B=(0,0,0)\)、\(A=(2,0,0)\)、\(C=\left(-\dfrac12,\dfrac{\sqrt3}{2},0\right)\). 由直三棱柱的高为 \(1\)，得 \(B_1=(0,0,1)\)、\(C_1=\left(-\frac12,\frac{\sqrt3}{2},1\right)\).
取两异面直线的方向向量
\(\overrightarrow{AB_1}=(-2,0,1)\)、\(\overrightarrow{BC_1}=\left(-\frac12,\frac{\sqrt3}{2},1\right)\).
它们的数量积为 \(2\)，长度分别为 \(\sqrt5\) 和 \(\sqrt2\). 因而所成角的余弦值为
\[
\frac{|\overrightarrow{AB_1}\cdot\overrightarrow{BC_1}|}{|\overrightarrow{AB_1}|\,|\overrightarrow{BC_1}|}
=\frac{2}{\sqrt5\sqrt2}=\frac{\sqrt{10}}5
\]
故选 C.
\end{solution}

\begin{problem}
若 \(x=-2\) 是函数 \(f(x)=(x^2+ax-1)\e^{x-1}\) 的极值点，则 \(f(x)\) 的极小值为（\quad）

\choices
    {\(-1\)}
    {\(-2\e^{-3}\)}
    {\(5\e^{-3}\)}
    {\(1\)}
\end{problem}

\begin{answer}
A
\end{answer}

\begin{solution}
求导得
\(f'(x)=\e^{x-1}\bigl[x^2+(a+2)x+a-1\bigr]\).
由于 \(\e^{x-1}>0\)，且 \(-2\) 是极值点，故
\(4-2(a+2)+a-1=0\)，
解得 \(a=-1\). 此时
\(f'(x)=\e^{x-1}(x+2)(x-1)\).
因此 \(f'(x)\) 在 \(x=1\) 处由负变正，\(x=1\) 是极小值点. 代入 \(f(x)=(x^2-x-1)\e^{x-1}\)，得极小值
\(f(1)=(1-1-1)\e^0=-1\).
故选 A.
\end{solution}

\begin{problem}
已知 \(\triangle ABC\) 是边长为 \(2\) 的等边三角形，\(P\) 为平面 \(ABC\) 内一点，则 \(\overrightarrow{PA}\cdot(\overrightarrow{PB}+\overrightarrow{PC})\) 的最小值是（\quad）

\choices
    {\(-2\)}
    {\(-\dfrac32\)}
    {\(-\dfrac43\)}
    {\(-1\)}
\end{problem}

\begin{answer}
B
\end{answer}

\begin{solution}
设 \(M\) 为 \(BC\) 的中点，则 \(\overrightarrow{PB}+\overrightarrow{PC}=2\overrightarrow{PM}\). 令 \(\bs{u}=\overrightarrow{PA}\)，\(\bs{v}=\overrightarrow{AM}\)，则 \(\overrightarrow{PM}=\bs{u}+\bs{v}\)，所以
\[
\overrightarrow{PA}\cdot(\overrightarrow{PB}+\overrightarrow{PC})
=2\bs{u}\cdot(\bs{u}+\bs{v})
=2\left|\bs{u}+\frac12\bs{v}\right|^2-\frac12|\bs{v}|^2
\]
等边三角形边长为 \(2\)，故中线 \(AM=\sqrt3\)，即 \(|\bs{v}|=\sqrt3\). 因而上式的最小值为
\(-\frac12(\sqrt3)^2=-\frac32\)，
且当 \(P\) 为 \(AM\) 的中点时可以取得. 故选 B.
\end{solution}

\section{填空题}

\begin{problem}
一批产品的二等品率为 \(0.02\)，从这批产品中每次随机取一件，有放回地抽取 \(100\text{ 次}\). \(X\) 表示抽到的二等品件数，则 \(D(X)=\)\fillinblank{}
\end{problem}

\begin{answer}
\(1.96\)
\end{answer}

\begin{solution}
每次抽取是否抽到二等品相互独立，且抽到二等品的概率为 \(p=0.02\). 因而 \(X\sim\mathrm{Bin}(100,0.02)\)，其方差为
\(D(X)=np(1-p)=100\times0.02\times0.98=1.96\).
所以填 \(1.96\).
\end{solution}

\begin{problem}
函数 \(f(x)=\sin^2x+\sqrt3\cos x-\dfrac34\)（\(x\in\left[0,\dfrac{\pi}{2}\right]\)）的最大值是\fillinblank{}
\end{problem}

\begin{answer}
1
\end{answer}

\begin{solution}
令 \(t=\cos x\)，则 \(0\le t\le1\)，且
\(f(x)=1-t^2+\sqrt3t-\frac34=1-\left(t-\frac{\sqrt3}{2}\right)^2\).
因为 \(\dfrac{\sqrt3}{2}\in[0,1]\)，所以当 \(t=\dfrac{\sqrt3}{2}\) 时，\(f(x)\) 取得最大值 \(1\).
\end{solution}

\begin{problem}
等差数列 \(\{a_n\}\) 的前 \(n\) 项和为 \(S_n\)，\(a_3=3\)，\(S_4=10\)，则 \(\displaystyle \sum_{k=1}^n \frac1{S_k}=\)\fillinblank{}
\end{problem}

\begin{answer}
\(\dfrac{2n}{n+1}\)
\end{answer}

\begin{solution}
设等差数列首项为 \(a_1\)，公差为 \(d\). 由 \(a_3=3\) 和 \(S_4=10\)，得
\(a_1+2d=3\)、\(4a_1+6d=10\).
解得 \(a_1=1\)，\(d=1\). 所以
\(S_k=\frac{k(a_1+a_k)}2=\frac{k(k+1)}2\)，从而
\[
\sum_{k=1}^n\frac1{S_k}=\sum_{k=1}^n\frac{2}{k(k+1)}
=2\sum_{k=1}^n\left(\frac1k-\frac1{k+1}\right)
=2\left(1-\frac1{n+1}\right)=\frac{2n}{n+1}
\]
所以填 \(\dfrac{2n}{n+1}\).
\end{solution}

\begin{problem}
已知 \(F\) 是抛物线 \(C:y^2=8x\) 的焦点，\(M\) 是 \(C\) 上一点，\(FM\) 的延长线交 \(y\) 轴于点 \(N\). 若 \(M\) 为 \(FN\) 的中点，则 \(|FN|=\)\fillinblank{}
\end{problem}

\begin{answer}
\(6\)
\end{answer}

\begin{solution}
由 \(y^2=8x\) 与标准方程 \(y^2=2px\) 比较，得 \(p=4\)，所以焦点为 \(F=(2,0)\). 设抛物线上一点 \(M\) 为
\(M=(2t^2,4t)\).
因为 \(M\) 是 \(FN\) 的中点，所以 \(N=2M-F=(4t^2-2,8t)\). 又 \(N\) 在 \(y\) 轴上，故 \(4t^2-2=0\)，即 \(t^2=\dfrac12\). 因此
\(|FN|=\sqrt{(0-2)^2+(8t)^2}=\sqrt{4+32}=6\).
所以填 \(6\).
\end{solution}

\section{解答题}

\begin{problem}
\(\triangle ABC\) 的内角 \(A\)，\(B\)，\(C\) 的对边分别为 \(a\)，\(b\)，\(c\)，已知 \(\sin(A+C)=8\sin^2\dfrac B2\).

\begin{enumerate}
\item 求 \(\cos B\)；
\item 若 \(a+c=6\)，\(\triangle ABC\) 的面积为 \(2\)，求 \(b\).
\end{enumerate}
\end{problem}

\begin{solution}
\begin{enumerate}
\item 因为 \(A+B+C=\pi\)，所以 \(A+C=\pi-B\)，从而 \(\sin B=\sin(A+C)=8\sin^2\frac B2\).
又因为 \(0<B<\pi\)，所以 \(\sin\dfrac B2>0\). 利用半角公式，得
\(2\sin\frac B2\cos\frac B2=8\sin^2\frac B2\).
因此
\(\cos\frac B2=4\sin\frac B2\)、\(\tan\frac B2=\frac14\).
于是
\[
\cos B=\frac{1-\tan^2\dfrac B2}{1+\tan^2\dfrac B2}
=\frac{1-\dfrac1{16}}{1+\dfrac1{16}}=\frac{15}{17}
\]

\item 由上问和 \(0<B<\pi\)，有 \(\sin B=\frac{2\tan\dfrac B2}{1+\tan^2\dfrac B2}=\frac8{17}\).
由三角形面积公式及 \(S_{\triangle ABC}=2\)，得
\(2=\frac12ac\sin B=\frac12ac\times\frac8{17}\)，所以 \(ac=\frac{17}{2}\).
再由余弦定理和 \(a+c=6\)，得
\[
b^2=a^2+c^2-2ac\cos B=(a+c)^2-2ac(1+\cos B)
=36-17\left(1+\frac{15}{17}\right)=4
\]
由于 \(b>0\)，所以 \(b=2\).
\end{enumerate}
\end{solution}

\begin{problem}
海水养殖场进行某水产品的新、旧网箱养殖方法的产量对比，收获时各随机抽取了 \(100\text{ 个}\)网箱，测量各箱水产品的产量（单位：\(\mathrm{kg}\)），其频率分布直方图如图：

\begin{enumerate}
\item 设两种养殖方法的箱产量相互独立，记 \(A\) 表示事件“旧养殖法的箱产量低于 \(50 \mathrm{~kg}\)，新养殖法的箱产量不低于 \(50 \mathrm{~kg}\)”，估计 \(A\) 的概率；
\item 填写下面列联表，并根据列联表判断是否有 \(99\%\) 的把握认为箱产量与养殖方法有关：

\begin{center}
\begin{tabular}{c|cc}
& 箱产量 \(<50 \mathrm{~kg}\) & 箱产量 \(\ge 50 \mathrm{~kg}\)\\
\hline
旧养殖法 & &\\
新养殖法 & &
\end{tabular}
\end{center}

\item 根据箱产量的频率分布直方图，求新养殖法箱产量的中位数的估计值（精确到 \(0.01\)）.
\end{enumerate}

\bitmapfigure[max width=0.92\linewidth]{img/2017/national_paper_2_science/q18_fig1.png}

附：

\begin{center}
\begin{tabular}{c|ccc}
\(P(K^2\ge k)\) & 0.050 & 0.010 & 0.001\\
\hline
k & 3.841 & 6.635 & 10.828
\end{tabular}
\end{center}

另参考公式 \(K^2=\dfrac{n(ad-bc)^2}{(a+b)(c+d)(a+c)(b+d)}\).

\end{problem}

\begin{solution}
\begin{enumerate}
\item 记事件 \(B\) 为“旧养殖法的箱产量低于 \(50\mathrm{~kg}\)”，事件 \(C\) 为“新养殖法的箱产量不低于 \(50\mathrm{~kg}\)”. 由旧养殖法直方图可得
\(P(B)=(0.012+0.014+0.024+0.034+0.040)\times5=0.62\).
由新养殖法直方图可得
\(P(C)=(0.068+0.046+0.010+0.008)\times5=0.66\).
因为两种养殖方法的箱产量相互独立，所以
\(P(A)=P(B)P(C)=0.62\times0.66=0.4092\).

\item 由两种养殖方法的频率分布直方图，得到列联表

\begin{center}
\begin{tabular}{c|cc}
& 箱产量 \(<50\mathrm{~kg}\) & 箱产量 \(\ge 50\mathrm{~kg}\)\\
\hline
旧养殖法 & 62 & 38\\
新养殖法 & 34 & 66
\end{tabular}
\end{center}

代入题给公式，得
\[
K^2=\frac{200(62\times66-38\times34)^2}{100\times100\times96\times104}\approx15.705
\]
由于 \(15.705>6.635\)，所以有 \(99\%\) 的把握认为箱产量与养殖方法有关.

\item 新养殖法箱产量低于 \(50\mathrm{~kg}\) 的频率为
\((0.004+0.020+0.044)\times5=0.34<0.5\)，
低于 \(55\mathrm{~kg}\) 的频率为
\((0.004+0.020+0.044+0.068)\times5=0.68>0.5\).
因此中位数的估计值在区间 \([50,55)\) 内. 设该估计值为 \(m\)，按组内频率均匀分布估计，有
\(0.34+0.068(m-50)=0.5\).
从而
\(m=50+\frac{0.5-0.34}{0.068}\approx52.35\mathrm{~kg}\).
所以新养殖法箱产量的中位数估计值为 \(52.35\mathrm{~kg}\).
\end{enumerate}
\end{solution}

\begin{problem}
如图，四棱锥 \(P-ABCD\) 中，侧面 \(PAD\) 为等边三角形且垂直于底面 \(ABCD\)，\(AB=BC=\dfrac12AD\)，\(\angle BAD=\angle ABC=90^\circ\)，\(E\) 是 \(PD\) 的中点.

\begin{enumerate}
\item 证明：直线 \(CE\parallel\) 平面 \(PAB\)；
\item 点 \(M\) 在棱 \(PC\) 上，且直线 \(BM\) 与底面 \(ABCD\) 所成角为 \(45^\circ\)，求二面角 \(M-AB-D\) 的余弦值.
\end{enumerate}

\bitmapfigure[max width=0.34\linewidth]{img/2017/national_paper_2_science/q19_fig1.png}
\end{problem}

\begin{solution}
\begin{enumerate}
\item 取 \(F\) 为 \(PA\) 的中点，连接 \(EF\) 和 \(BF\). 在三角形 \(PAD\) 中，\(E\)，\(F\) 分别为 \(PD\)，\(PA\) 的中点，所以
\(EF\parallel AD\)、\(EF=\frac12AD\).
由 \(\angle BAD=\angle ABC=90^\circ\) 及底面共面性，得 \(AD\parallel BC\). 又因为 \(BC=\dfrac12AD\)，所以 \(EF\parallel BC\) 且 \(EF=BC\). 因此四边形 \(BCEF\) 是平行四边形，从而 \(CE\parallel BF\).

因为 \(BF\subset\) 平面 \(PAB\)，且 \(CE\) 不在平面 \(PAB\) 内，所以 \(CE\parallel\) 平面 \(PAB\).

\item 令 \(AB=BC=1\)，则 \(AD=2\). 建立空间直角坐标系，使底面 \(ABCD\) 在 \(z=0\) 内，取 \(A(0,0,0)\)、\(B(1,0,0)\)、\(C(1,1,0)\)、\(D(0,2,0)\).
由于侧面 \(PAD\) 是边长为 \(2\) 的等边三角形且垂直于底面，可取
\(P(0,1,\sqrt3)\).
设 \(M=P+\lambda(C-P)\)，其中 \(0\le\lambda\le1\)，并令 \(t=1-\lambda\). 则
\(M(1-t,1,\sqrt3t)\)、\(\overrightarrow{BM}=(-t,1,\sqrt3t)\).
直线 \(BM\) 与底面所成角为 \(45^\circ\)，故
\[
\frac{\sqrt3t}{\sqrt{t^2+1+3t^2}}=\sin45^\circ=\frac{\sqrt2}{2}
\]
由于 \(0\le t\le1\)，解得 \(t=\dfrac1{\sqrt2}\). 因而点 \(M\) 的竖直坐标为
\(z_M=\frac{\sqrt6}{2}\).

平面 \(MAB\) 的一个法向量为 \(\bs{n}_1=\overrightarrow{AB}\times\overrightarrow{AM}=\left(0,-\frac{\sqrt6}{2},1\right)\)，平面 \(DAB\) 的一个法向量为 \(\bs{n}_2=(0,0,1)\). 设二面角 \(M-AB-D\) 为 \(\alpha\)，则
\[
\cos\alpha=\frac{|\bs{n}_1\cdot\bs{n}_2|}{|\bs{n}_1||\bs{n}_2|}
=\frac1{\sqrt{1+\dfrac32}}=\frac{\sqrt{10}}5
\]
\end{enumerate}
\end{solution}

\begin{problem}
设 \(O\) 为坐标原点，动点 \(M\) 在椭圆 \(C:\dfrac{x^2}{2}+y^2=1\) 上，过 \(M\) 作 \(x\) 轴的垂线，垂足为 \(N\)，点 \(P\) 满足 \(\overrightarrow{NP}=\sqrt2\,\overrightarrow{NM}\).

\begin{enumerate}
\item 求点 \(P\) 的轨迹方程；
\item 设点 \(Q\) 在直线 \(x=-3\) 上，且 \(\overrightarrow{OP}\cdot \overrightarrow{PQ}=1\). 证明：过点 \(P\) 且垂直于 \(OQ\) 的直线 \(l\) 过 \(C\) 的左焦点 \(F\).
\end{enumerate}
\end{problem}

\begin{solution}
\begin{enumerate}
\item 设 \(M(x_0,y_0)\)，则 \(N(x_0,0)\). 设 \(P(x,y)\)，由 \(\overrightarrow{NM}=(0,y_0)\)、\(\overrightarrow{NP}=(x-x_0,y)\)，以及 \(\overrightarrow{NP}=\sqrt2\,\overrightarrow{NM}\)，得到 \(x=x_0\)、\(y=\sqrt2y_0\).
因为 \(M\) 在椭圆 \(C\) 上，所以
\(\frac{x_0^2}{2}+y_0^2=1\).
代入 \(x=x_0\)，\(y=\sqrt2y_0\)，得
\(x^2+y^2=2\).
反之，圆 \(x^2+y^2=2\) 上任一点 \(P(x,y)\) 均可取 \(M\left(x,\dfrac y{\sqrt2}\right)\)，故点 \(P\) 的轨迹为
\(x^2+y^2=2\).

\item 椭圆 \(C\) 的焦半距为 \(\sqrt{2-1}=1\)，所以左焦点为 \(F(-1,0)\). 设 \(Q(-3,t)\)，\(P(x,y)\). 由上问 \(x^2+y^2=2\)，且
\[
1=\overrightarrow{OP}\cdot\overrightarrow{PQ}
=(x,y)\cdot(-3-x,t-y)
=-3x-x^2+ty-y^2
=-3x-2+ty
\]
因此 \(ty=3(x+1)\).
又
\[
\overrightarrow{FP}\cdot\overrightarrow{OQ}=(x+1,y)\cdot(-3,t)=-3(x+1)+ty=0
\]
所以 \(FP\perp OQ\). 过 \(P\) 且垂直于 \(OQ\) 的直线经过 \(F\)，即该直线 \(l\) 过椭圆 \(C\) 的左焦点 \(F\).
\end{enumerate}
\end{solution}

\begin{problem}
已知函数 \(f(x)=ax^2-ax-x\ln x\)，且 \(f(x)\ge 0\).

\begin{enumerate}
\item 求 \(a\)；
\item 证明：\(f(x)\) 存在唯一的极大值点 \(x_0\)，且 \(\e^{-2}<f(x_0)<2^{-2}\).
\end{enumerate}
\end{problem}

\begin{solution}
\begin{enumerate}
\item 函数的定义域为 \(x>0\). 将已知不等式写成
\(f(x)=x\bigl(a(x-1)-\ln x\bigr)\ge0\).
令 \(g(x)=a(x-1)-\ln x\)，则 \(g(x)\ge0\) 且 \(g(1)=0\). 因为 \(1\) 是定义域内点，\(g\) 在 \(x=1\) 处取得最小值，所以
\(g'(1)=a-1=0\).
故 \(a=1\).

\item 由 \(a=1\)，有 \(f(x)=x^2-x-x\ln x\)，\(f'(x)=2x-2-\ln x\).
令 \(h(x)=f'(x)\)，则
\(h'(x)=2-\frac1x\).
 因此 \(h\) 在 \(\left(0,\dfrac12\right)\) 上单调递减，在 \(\left(\dfrac12,+\infty\right)\) 上单调递增. 又
\(\displaystyle\lim_{x\to0^+}h(x)=+\infty\)、\(h\left(\frac12\right)=\ln2-1<0\)、\(h(1)=0\).
所以 \(h\) 在 \(\left(0,\dfrac12\right)\) 内有唯一零点，记为 \(x_0\)，而在 \(\left(\dfrac12,+\infty\right)\) 内只有零点 \(1\). 由 \(h\) 的符号变化，\(f\) 在 \((0,x_0)\) 上递增，在 \((x_0,1)\) 上递减，在 \((1,+\infty)\) 上递增，故 \(x_0\) 是唯一的极大值点.

由 \(h(x_0)=0\)，得
\(\ln x_0=2x_0-2\).
于是
\(f(x_0)=x_0^2-x_0-x_0\ln x_0=x_0(1-x_0)\).
因为 \(0<x_0<\dfrac12\)，所以
\(f(x_0)<\frac14\).
另一方面，令 \(\varphi(x)=\ln(1-x)+2x\)，则在 \(0<x<\dfrac12\) 时
\(\varphi'(x)=\frac{1-2x}{1-x}>0\)，且 \(\displaystyle\lim_{x\to0^+}\varphi(x)=0\).
因此 \(\ln(1-x)>-2x\)，即 \(1-x>\e^{-2x}\). 取 \(x=x_0\)，并利用 \(x_0\e^{-2x_0}=\e^{-2}\)，得
\(f(x_0)=x_0(1-x_0)>x_0\e^{-2x_0}=\e^{-2}\).
综上，
\(\e^{-2}<f(x_0)<2^{-2}\).
\end{enumerate}
\end{solution}

\section{选考题：解答题}

\begin{problem}
在直角坐标系 \(xOy\) 中，以坐标原点为极点，\(x\) 轴的正半轴为极轴建立极坐标系，曲线 \(C_1\) 的极坐标方程为 \(\rho\cos\theta=4\).

\begin{enumerate}
\item \(M\) 为曲线 \(C_1\) 上的动点，点 \(P\) 在线段 \(OM\) 上，且满足 \(|OM|\cdot|OP|=16\)，求点 \(P\) 的轨迹 \(C_2\) 的直角坐标方程；
\item 设点 \(A\) 的极坐标为 \(\left(2,\dfrac{\pi}{3}\right)\)，点 \(B\) 在曲线 \(C_2\) 上，求 \(\triangle OAB\) 面积的最大值.
\end{enumerate}
\end{problem}

\begin{solution}
\begin{enumerate}
\item 设 \(M\) 的极坐标为 \((\rho,\theta)\)，点 \(P\) 的极径为 \(r\). 因为 \(P\) 在线段 \(OM\) 上，且
\(\rho r=|OM||OP|=16\).
又 \(M\in C_1\)，所以 \(\rho\cos\theta=4\). 从而
\(r=\frac{16}{\rho}=4\cos\theta\).
由于 \(r>0\)，有 \(\cos\theta>0\). 将 \(x=r\cos\theta\)，\(y=r\sin\theta\) 代入，得
\(r^2=4r\cos\theta\)、\(x^2+y^2=4x\).
因此点 \(P\) 的轨迹为 \((x-2)^2+y^2=4\)，\(x>0\).
反过来，设圆上任意一点 \(P(x,y)\) 满足 \(x>0\)，令 \(r=\sqrt{x^2+y^2}\)，并取 \(\theta\) 使 \(\cos\theta=x/r\)、\(\sin\theta=y/r\). 由 \(x^2+y^2=4x\) 得 \(r=4\cos\theta\)，从而 \(\cos\theta>0\). 再取 \(\rho=4/\cos\theta\)，则点 \(M\) 的极坐标 \((\rho,\theta)\) 满足 \(\rho\cos\theta=4\)、\(\rho r=16\)，且由 \(P\) 与 \(M\) 同向可知 \(P\) 满足题设关系. 故这就是轨迹方程.

\item 点 \(A\) 的直角坐标为
\(A\left(2\cos\frac\pi3,2\sin\frac\pi3\right)=(1,\sqrt3)\).
设 \(B(x,y)\in C_2\). 则
\[
S_{\triangle OAB}=\frac12\left|\det(\overrightarrow{OA},\overrightarrow{OB})\right|
=\frac12|y-\sqrt3x|
\]
将圆参数化为
\(x=2+2\cos t\)、\(y=2\sin t\).
可得
\(y-\sqrt3x=2\bigl(\sin t-\sqrt3\cos t-\sqrt3\bigr)\).
由
\(-2\le\sin t-\sqrt3\cos t\le2\)，
可知
\(-4-2\sqrt3\le y-\sqrt3x\le4-2\sqrt3\).
因此 \(|y-\sqrt3x|\) 的最大值为 \(4+2\sqrt3\)，例如在 \(t=-\dfrac\pi6\) 时取得，且此时 \(x=2+\sqrt3>0\). 所以
\[
S_{\triangle OAB,\max}=\frac12(4+2\sqrt3)=2+\sqrt3
\]
\end{enumerate}
\end{solution}

\begin{problem}
已知 \(a>0\)，\(b>0\)，\(a^3+b^3=2\). 证明：

\begin{enumerate}
\item \((a+b)(a^5+b^5)\ge 4\)；
\item \(a+b\le 2\).
\end{enumerate}
\end{problem}

\begin{solution}
\begin{enumerate}
\item 由柯西不等式，得
\[
(a+b)(a^5+b^5)\ge\bigl(\sqrt a\,a^{\frac52}+\sqrt b\,b^{\frac52}\bigr)^2
=(a^3+b^3)^2=4
\]

\item 由霍尔德不等式，得
\[
(a+b)^3=(a\cdot1+b\cdot1)^3
\le(a^3+b^3)(1^3+1^3)^2=2\times4=8
\]
因为 \(a+b>0\)，所以 \(a+b\le2\).
\end{enumerate}
\end{solution}
